\DocumentMetadata{
  pdfversion=2.0,
  pdfstandard=ua-2,
  lang=en-US,
  tagging=on,
  tagging-setup={math/setup=mathml-SE,tabsorder=structure}
}
\documentclass[11pt,letterpaper]{article}
\usepackage[margin=0.68in,includefoot]{geometry}
\usepackage{newcomputermodern}
\usepackage{amsmath,amscd,mathtools}
\usepackage{tikz,adjustbox}
\providecommand{\square}{\mdlgwhtsquare}
\usepackage[hidelinks]{hyperref}
\hypersetup{
  pdftitle={Algebra PhD Exam, January 2019},
  pdfauthor={University of Florida Department of Mathematics},
  pdfsubject={Graduate mathematics examination},
  pdfdisplaydoctitle=true,
  bookmarksopen=true
}
\setcounter{secnumdepth}{0}
\setlength{\parindent}{0pt}
\setlength{\parskip}{0.28em}
\setlength{\emergencystretch}{3em}
\setlength{\leftmargini}{2.15em}
\setlength{\leftmarginii}{2.65em}
\setlength{\leftmarginiii}{3em}
\renewcommand{\labelenumi}{\textbf{\theenumi.}}
\pagestyle{empty}
\raggedbottom
\allowdisplaybreaks
\newenvironment{examproblems}[1]{%
  \begin{enumerate}%
  \setcounter{enumi}{#1}%
  \setlength{\topsep}{0.35em}%
  \setlength{\partopsep}{0pt}%
  \setlength{\itemsep}{0.48em}%
  \setlength{\parsep}{0.18em}%
}{\end{enumerate}}
\newenvironment{examsubparts}{%
  \begin{enumerate}%
  \setlength{\topsep}{0.18em}%
  \setlength{\partopsep}{0pt}%
  \setlength{\itemsep}{0.16em}%
  \setlength{\parsep}{0.1em}%
}{\end{enumerate}}
\newenvironment{examinstructions}{%
  \par\begingroup\itshape
}{%
  \par\endgroup\vspace{0.18em}
}
\makeatletter
\renewcommand\section{\@startsection{section}{1}{0pt}%
  {-1.25ex plus -.25ex minus -.1ex}%
  {0.55ex plus .1ex}%
  {\normalfont\large\bfseries}}
\renewcommand{\@maketitle}{%
  \newpage\null\vskip -1em%
  {\footnotesize\bfseries
    \href{https://www.ufl.edu/}{University of Florida}, \href{https://math.ufl.edu/}{Department of Mathematics}\par}%
  \vskip -0.5em%
  \tagstructbegin{tag=Title}%
    {\LARGE\bfseries\href{https://gma.math.ufl.edu/exams/algebra-phd/}{Algebra PhD Exam}, January 2019\par}%
  \tagstructend%
  \vskip 0.25em%
  \tagmcbegin{artifact}%
    \hrule height 1pt%
  \tagmcend%
  \vskip 1em%
}
\makeatother
\title{Algebra PhD Exam, January 2019}
\author{}
\date{}
\begin{document}
\maketitle
\thispagestyle{empty}
\begin{examinstructions}
Answer seven problems. Write your answers clearly in complete English sentences. You may quote results (within reason) as long as you state them clearly.
\end{examinstructions}
\begin{examproblems}{0}
\item
Determine the Galois group of the polynomial 
\[
f(x)=x^4-14x^2+9\in\mathbb{Q}[x].
\]
\item
This problem has three parts.
\begin{examsubparts}
\item[\textnormal{(a)}]
Give the definition of a separable field extension.
\item[\textnormal{(b)}]
Show that a finite separable extension is contained in a finite Galois extension.
\item[\textnormal{(c)}]
Prove the Primitive Element Theorem: If~\(E\) is a finite separable extension of~\(F\), then \(E=F(\alpha)\) for some \(\alpha\in E\).
\end{examsubparts}
\item
A pointed set is a pair \((S,x)\) with~\(S\) a set and \(x\in S\). A morphism of pointed sets \((S,x)\to(S',x')\) is a function \(f:S\to S'\) such that \(f(x)=x'\).
\begin{examsubparts}
\item[\textnormal{(a)}]
Show that pointed sets and their morphisms form a category.
\item[\textnormal{(b)}]
Show that coproducts of arbitrarily indexed families of objects exist in this category.
\end{examsubparts}
\item
Let~\(F\) be the free group on~\(n\) generators, \(n\geq 1\).
\begin{examsubparts}
\item[\textnormal{(a)}]
Prove that the automorphism group of~\(F\) contains a subgroup isomorphic to the symmetric group \(S_n\).
\item[\textnormal{(b)}]
Give an example of an automorphism of~\(F\) that does not belong to this subgroup.
\end{examsubparts}
\item
Define the term injective module for a ring~\(R\) with 1. Prove Baer’s criterion: A unitary left module~\(J\) is injective if and only if for every left ideal~\(L\) of~\(R\), any~\(R\)-module homomorphism \(L\to J\) can be extended to an~\(R\)-module homomorphism \(R\to J\).
\item
Let~\(R\) be a ring and suppose that we have a short exact sequence 
\[
0\to A\xrightarrow{f}B\xrightarrow{g}C\to 0
\]
 of left~\(R\)-modules. Let~\(M\) be a right~\(R\)-module. Show that the induced sequence 
\[
M\otimes_R A\xrightarrow{1\otimes f}M\otimes_R B\xrightarrow{1\otimes g}M\otimes_R C\to 0
\]
 is exact.
\item
Let~\(R\) be a commutative ring with \(1\neq 0\) and let~\(S\) be a multiplicative set containing 1.
\begin{examsubparts}
\item[\textnormal{(a)}]
State the universal mapping property of the localization \(S^{-1}R\).
\item[\textnormal{(b)}]
Show that if \(S=R\setminus P\), the complement of a prime ideal~\(P\), then \(S^{-1}R\) has a unique maximal ideal.
\end{examsubparts}
\item
Prove the following properties of \(R=\mathbb{Z}[\sqrt{10}]\).
\begin{examsubparts}
\item[\textnormal{(a)}]
\(R\) is noetherian.
\item[\textnormal{(b)}]
\(R\) is an integrally closed integral domain.
\item[\textnormal{(c)}]
Every nonzero prime ideal of~\(R\) is maximal.
\end{examsubparts}
\item
Let~\(J\) be an ideal in a commutative ring~\(R\) with 1. Assume that~\(J\) is contained in every maximal ideal. Prove that if~\(A\) is an~\(R\)-module that satisfies the ascending chain condition on submodules, then \(\bigcap_{n=1}^{\infty}J^nA=0\).
\item
Prove the Noether Normalization Lemma: Let~\(R\) be an integral domain that is a finitely generated extension ring of a field~\(K\). Let~\(F\) be the field of quotients of~\(R\) and let~\(r\) be the transcendence degree of~\(F\) over~\(K\). Then there exist a subset of~\(r\) elements \(t_1,\ldots,t_r\) which are algebraically independent over~\(K\) and such that~\(R\) is integral over \(K[t_1,\ldots,t_r]\).
\item
This problem has two parts.
\begin{examsubparts}
\item[\textnormal{(a)}]
Define the term left primitive ring and state Jacobson’s Density Theorem.
\item[\textnormal{(b)}]
Prove that if~\(R\) is a finite left primitive ring then \(R\cong M_n(F)\) for some finite field~\(F\) and some natural number~\(n\).
\end{examsubparts}
\end{examproblems}
\end{document}
